\section{AI 测试技术落地使用情况概述}

\subsection{测试对象与判定原则}

本作品的测试对象分为两类。第一类是确定性代码，包括表达式白名单、答案归一化、符号验算、检索排序、数据库迁移、接口参数和事件顺序。这类代码必须有明确输入与期望输出，失败时能够定位到具体规则。第二类是模型相关路径，包括视觉识别、卡片萃取、命题、讲解和批改。这类输出具有波动性，测试不能只比较整段文案，也不能用“看起来合理”作为通过标准。

因此，模型相关测试主要检查契约和系统行为：返回结构能否被收敛，非法字段能否被丢弃，流式片段能否无丢失拼接，模型失败时是否进入声明过的降级路径，真实 Provider 缺配置时是否明确报未就绪。真实联调用于确认外部协议与主路径可用，不与离线测试混为一谈。所有最终断言仍是确定性条件，模型本身不担任裁判。

测试设计遵循三条边界。其一，单元测试不访问真实模型，也不碰开发数据库。其二，评测数字必须说明样本构造方式，不能把样例通过率写成泛化能力。其三，安全结论必须实际执行攻击输入并检查副作用，不能只阅读代码后声称安全。

\begin{table}[H]
\centering\footnotesize
\caption{测试分层、工具与各自能够证明的事实}
\begin{tabularx}{\textwidth}{@{}L{2.4cm}L{3.0cm}Y L{3.5cm}@{}}
\toprule
\textbf{层级} & \textbf{工具与范围} & \textbf{主要断言} & \textbf{不能替代的检查} \\
\midrule
后端单元与接口 & pytest，26 个测试文件，403 项 & 规则结果、异常分支、HTTP 状态、数据库变化、事件顺序与配置回退 & 真实模型质量、完整浏览器交互 \\
前端纯函数 & Node 自带测试器，3 个文件，18 项 & 上传闸门、题数夹取、流式字段归一、批改清洗、回放结局合成 & 组件布局、浏览器兼容性 \\
类型与生产构建 & \env{tsc --noEmit} 与 Vite build & 类型错误、模块解析、静态产物可生成 & 运行时业务正确性 \\
端到端冒烟 & Playwright + Chromium，7 条主路径 & 页面按钮、数据加载、题目回环、讲解、批改、修订与回放真实连通 & 长稳、并发与真实模型波动 \\
专项构造评测 & 离线 Python 脚本与结果 JSON & 坏题拦截、好题误驳、结构检查、安全载荷、检索命中率 & 大规模真实题库上的统计结论 \\
真实 Provider 冒烟 & 智谱视觉 + DeepSeek 官方文本 & 外部协议、真实识别、真实驳回重生成与批改路径可走通 & 免费档以外模型能力、跨学科准确率 \\
\bottomrule
\end{tabularx}
\end{table}

\subsection{自动化回归结果与隔离方式}

当前质量基线为后端 403 项测试全绿，实测约 68 s；前端 18 项全绿；类型检查和生产构建通过；七条端到端路径在源码模式与打包版分别通过。后端数量相较早期技术文档中的 395 项有所增加，本文采用 v0.1.6 的现状口径。测试数量本身不是质量结论，重要的是是否覆盖曾经出错的边界以及修复后能否防止回归。

\begin{table}[H]
\centering\footnotesize
\caption{自动化回归结果摘要}
\begin{tabularx}{\textwidth}{@{}L{3.0cm}L{3.2cm}Y@{}}
\toprule
\textbf{检查项} & \textbf{当前结果} & \textbf{证据范围与复现入口} \\
\midrule
后端 pytest & 403 项，26 个文件，全绿 & \file{python -m pytest apps/server/tests -q}；覆盖验算、判分、检索、接口、事件、文档与安全 \\
前端单元测试 & 18 项，全部通过 & \file{cd apps/web && npm test}；不引入额外测试框架，直接检查纯函数 \\
前端构建 & 类型检查与 Vite 构建通过 & \file{npm run build} 先执行 \env{tsc --noEmit}，再生成静态资源 \\
Playwright 冒烟 & 7/7；源码版与打包版各一轮 & \file{python scripts/e2e_smoke.py --shots ./shots}；失败时保存截图 \\
桌面运行时检查 & 健康检查、静态页面与 CSP 可用 & 使用独立临时数据目录；退出后检查后端进程，不写入开发库 \\
\bottomrule
\end{tabularx}
\end{table}

后端 \file{conftest.py} 自动把数据库和上传目录指向 pytest 临时路径，并强制 mock，开发机即使配置真实 Key 也不会在测试中联网。E2E 在启动前检查端口，被占即退出码 2，避免写入已有开发服务。

断言同时检查过程与结果：第 1 轮第 4 题必须失败、第 2 轮全过且命题恰好两次；事件裁剪不得丢终止状态；同值修订不写历史；检索故障必须退回 2-gram；判分同时保留应判对与应判错样例。

\subsection{专项评测：数字、图形与口径}

专项评测脚本将逐题结果写入 \file{scripts/eval-results/}，摘要数字由结果计算而非手工填入。验算集从三套教学样例取 15 道题，将每题的 \env{expect} 偏移 1 构造不自洽答案；同一批未修改题用于误驳检查。结构集包含空题干、空答案和坏选项等五种损坏形态。检索集由 20 个问句与 10 张卡片构成，包含同义改写和英文问句，判定正确卡片是否进入前三名。

\begin{table}[H]
\centering\footnotesize
\caption{专项评测结果与统计口径}
\begin{tabularx}{\textwidth}{@{}L{2.8cm}L{2.7cm}Y L{3.4cm}@{}}
\toprule
\textbf{指标} & \textbf{结果} & \textbf{构造与计算方式} & \textbf{结论边界} \\
\midrule
不自洽题拦截率 & 15/15，100\% & 15 道正确题的期望值偏移 1，逐题调用当前质检逻辑 & 证明构造错误可拦截，不等于真实模型错误率 \\
正确题误驳率 & 0/15，0\% & 同一批原始正确题全部通过 & 样本量小，不覆盖复杂数学式 \\
结构损坏拦截率 & 5/5，100\% & 无验算账单且缺题干、答案或选项结构不合法 & 只证明形状检查有效 \\
攻击载荷拦截 & 3/3，100\% & 摘要脚本中的注入与指数塔输入 & 与更完整的安全回归和独立六载荷复验互相补充 \\
2-gram top-3 & 16/20，80\% & 20 问句对 10 张卡，正确卡片进入前三记命中 & 小语料，不能外推到大库 \\
混合检索 top-3 & 19/20，95\% & 本地 BGE-M3 q8\_0，2-gram 与向量按 RRF 融合，阈值 0.35 & 依赖嵌入端点；唯一未命中为英文跨语问句 \\
\bottomrule
\end{tabularx}
\end{table}

\begin{figure}[H]
\centering
\begin{tikzpicture}
\begin{axis}[
  width=0.88\textwidth,height=6.3cm,
  ybar,bar width=22pt,
  ymin=0,ymax=110,
  ylabel={通过率或命中率（\%）},
  symbolic x coords={验算拦截,结构拦截,2-gram 检索,混合检索},
  xtick=data,
  x tick label style={font=\small},
  y tick label style={font=\small},
  ylabel style={font=\small},
  nodes near coords={\pgfmathprintnumber{\pgfplotspointmeta}\%},
  nodes near coords style={font=\small,color=ink},
  axis lines=left,
  axis line style={line},
  tick style={line},
  ymajorgrids=true,
  grid style={line!55},
  every axis plot/.append style={fill=inklight,draw=ink}
]
\addplot coordinates {(验算拦截,100) (结构拦截,100) (2-gram 检索,80) (混合检索,95)};
\end{axis}
\end{tikzpicture}
\caption{构造集上的四项结果。前两项是缺陷拦截率，后两项是检索 top-3 命中率，统计对象不同，不应相加或解释为同一准确率。}
\label{fig:eval-bars}
\end{figure}

混合检索唯一未命中的问句是“how to find the vertex of a parabola”。该样本保留在结果文件中，没有因为影响指标而删除。95\% 只表示这个小集合上比基线多命中 3 个问句，不能写成检索能力提升了固定百分点后便可推广。下一阶段若扩展到 100--500 张卡片，需要沿用同一命中定义，并同时记录延迟、嵌入成本和失败回退次数。

\subsection{安全复验：从代码审阅转向攻击实测}

独立评审用 \env{__import__} 与文件写入表达式确认早期通用符号解析会实际执行代码，该问题被列为阻塞项。修复统一经过 \env{safe_math}，并补充嵌套幂限制。复验同时检查返回值和目录副作用，后续回归扩展到属性访问、列表推导、系统调用和非数值代入，也保留合法公式防止白名单过严。

\begin{table}[H]
\centering\footnotesize
\caption{独立安全复验的六类载荷与判定}
\begin{tabularx}{\textwidth}{@{}L{1.1cm}L{4.6cm}Y L{2.4cm}@{}}
\toprule
\textbf{编号} & \textbf{载荷代表} & \textbf{攻击目的与检查点} & \textbf{修复后结果} \\
\midrule
P1 & \env{__import__('os').getcwd()} & 动态导入并调用系统模块；检查是否执行及是否拒绝 & 拒绝，无副作用 \\
P2 & \file{open('pwned.txt','w').write('x')} & 通过表达式写文件；除返回值外检查目标文件 & 拒绝，文件不存在 \\
P3 & \env{eval('1+1')} & 尝试二次求值；验证内置函数不可达 & 拒绝 \\
P4 & \env{exec('x = 1')} & 尝试执行语句；验证表达式入口不暴露执行能力 & 拒绝 \\
P5 & \env{getattr(__builtins__, 'eval')('1')} & 通过内置命名空间间接取回危险函数 & 拒绝 \\
P6 & \file{().__class__.__bases__[0].__subclasses__()} & 通过对象属性链寻找可用类并逃逸 & 拒绝，无属性访问 \\
\bottomrule
\end{tabularx}
\end{table}

这次处置留下的测试原则是：解析不可信输入时，黑名单不够，必须把允许字符、允许名称、允许代入类型和允许语法结构写成正向集合。安全性也不能只看异常文本。文件写入载荷必须验证文件没有生成，资源消耗载荷必须验证在求值前快速返回。六载荷通过不意味着数学解析已经形式化证明安全，它表示已知攻击路径和后续回归样本均被覆盖。

\subsection{模型输出契约、故障注入与降级测试}

模型输出会包含代码围栏、解释性尾句、尾随逗号或 LaTeX 转义错误。测试针对这些实际形态分别构造输入。括号扫描必须忽略字符串内部的花括号，不能用贪婪正则把 JSON 后的解释一并截取；含 \env{$...$} 与单反斜杠字母串时，要修复可能被 JSON 当成控制字符的 LaTeX；完全不是 JSON 的输入仍应报错，不能被修复器变成另一个含义。

批改还检查字段收敛。分数被夹到 0--100，结论只允许三种枚举，步骤号越界时丢弃首个错误字段，字符串 \env{"false"} 不能被 Python 的真值规则当成真。模型结果缺字段时使用本地规则补齐，标准答案始终取数据库题目而不是模型自己重复的一份。解析失败时，事件带 \env{json-fallback}，界面显示本地规则生成的人读评语，不展示原始 JSON。

\begin{table}[H]
\centering\footnotesize
\caption{模型与外部依赖故障注入矩阵}
\begin{tabularx}{\textwidth}{@{}L{3.2cm}Y L{3.4cm}@{}}
\toprule
\textbf{注入条件} & \textbf{期望系统行为} & \textbf{验证方式} \\
\midrule
JSON 前后有解释或代码围栏 & 截取括号配平的对象，保留字符串内部花括号 & Provider 解析单元测试 \\
LaTeX 单反斜杠、尾随逗号 & 在限定条件下修复；无法修复则明确失败 & JSON 修复正反样例 \\
模型字段类型错误或缺字段 & 收敛、裁剪或用本地规则补齐，不破坏响应结构 & 批改与讲解收敛测试 \\
讲解参数为 NaN、字符串或反向区间 & 服务端不下发非法参数，前端只保留文字讲解 & 服务端与前端双层守卫 \\
流式中途失败 & 丢弃半截文本，回退完整结果，避免“半句加全文” & 中途异常注入与前端清洗 \\
SSE 断流或服务重启 & 先重放内存，再读落库事件，最后才重新 POST & 事件回放测试与 E2E \\
嵌入端点未配置、超时或协议异常 & 退回 legacy，回答与出处仍返回，事件说明回退 & MockTransport 与缓存故障测试 \\
事件数据库写入失败 & 当前实时事件继续推送，仅失去事后回放 & 发布函数异常注入 \\
\bottomrule
\end{tabularx}
\end{table}

\subsection{AI 在开发测试中的实际用法与意图控制}

AI 编码助手用于提出候选边界、生成模块草稿、补充测试和进行独立评审。开发者负责确定范围、审查关键逻辑、运行验收并决定是否采纳。每个核心会话围绕一个可验收问题展开，结束时应留下代码变化与对话快照。快照包含目标、关键提示词、输出摘要、人工判断和落地结果，其中否决记录用于说明为什么某些建议没有进入产品。

意图控制采用固定四段：目标、约束、验收、回退。目标必须能用一个结果判定，约束限定目录、依赖和禁止操作，验收给出命令与期望输出，回退要求新行为可配置或能恢复旧路径。上下文只提供相关文件和真实报错，不把整仓无筛选地交给模型。AI 不得自行引入新依赖，后端 8 个直接依赖和前端 3 个运行时依赖均受这一约束。

\begin{table}[H]
\centering\footnotesize
\caption{AI 协同开发中的采纳与否决实例}
\begin{tabularx}{\textwidth}{@{}L{2.8cm}L{4.4cm}Y@{}}
\toprule
\textbf{事项} & \textbf{AI 建议或输出} & \textbf{人工判断与最终验证} \\
\midrule
编排方案 & 使用 LangGraph 管理状态与检查点 & 否决。当前只有一个有限反馈回路，自研状态机更易于工期内调试；以流水线事件测试验收 \\
中文检索 & 按空白分词 & 否决。实测“为什么这里要配方”召回为 0；改 2-gram 后用固定问句回归 \\
表达式安全 & 过滤 \env{import}、\env{open} 等黑名单 & 否决。绕过形式过多；改字符、名称、代入和结构白名单，并做六载荷实测 \\
 混合检索 & 缺省启用向量并引入向量库 & 部分否决。保留可选 BGE-M3 + RRF，不引入原生向量库，缺省仍为 legacy \\
 \bottomrule
\end{tabularx}
\end{table}

中文召回为 0、mock 从不驳回和文档遗留缺陷都经过真实运行发现，再转为固定样例与回归测试。AI 提供草稿和候选用例，问题是否成立仍由执行结果决定。

\subsection{质量证据的诚实边界}

当前没有覆盖率统计、组件渲染测试和并发压测，E2E 也只跑 mock。真实联调覆盖智谱 \env{glm-4v-flash} 与 DeepSeek 官方 \env{deepseek-chat} 的主路径，不代表其他 Provider 均已回归；\env{student_method} 的真实样本验证仍不足。视觉免费档在复杂手写、表格和公式上的稳定性没有统计结论。
